\section{Appendix: Scriptural Date and Location}
\label{historical-appendix}
% Generated. Do not edit.
% Deterministic projection of research/chronology.toml.
% schema: 5
% document: liturgy/roman-rite/postconciliar/roman-missal-third-edition-en-us-2011/propers/general-calendar/pc-r07-all-saints-abc
% generated-by: tools/tpt proper-chronology annotations
\newcommand{\chronologyannotationclaim}[6]{#6}
\newcommand{\chronologyannotationreach}[3]{}
\newcommand{\chronologyannotationgroup}[3]{#3}
\newcommand{\chronologyannotation}[1]{%
  \ifcsname triptychchronologyannotation@#1\endcsname
    \csname triptychchronologyannotation@#1\endcsname
  \else
    \PackageError{triptych}{No chronology annotation for #1}{}%
  \fi}
% appointment-dependency: guidance/liturgy/postconciliar-propers-registry.md sha256=ffe8a6d8b7d79c652bdc49b3d5f6d4dbb32e3506abc7ab0310e745c2b7306ad4
% appointment-dependency: src/gpt/liturgy/roman-rite/postconciliar/roman-missal-third-edition-en-us-2011/propers/general-calendar/pc-r07-all-saints-abc/instance/manifest.md sha256=9ae6055704d1fe2aa3949d5e91e816673c113bc775b3a39b1bc145c4290b6630
% appointment-dependency: src/gpt/liturgy/roman-rite/postconciliar/roman-missal-third-edition-en-us-2011/propers/general-calendar/pc-r07-all-saints-abc/propers/verified.md sha256=574617c10630e751cb43a297cfec38b5c622631c3dd5ca4b8092d3f70c5441b5
% appointment-dependency: src/gpt/liturgy/roman-rite/postconciliar/roman-missal-third-edition-en-us-2011/propers/general-calendar/pc-r07-all-saints-abc/research/chronology-inputs.toml sha256=f30b32a467cb7d56012eddb3291e509def6fa3b786ee844bead09d3b8452e975
% appointment-dependency: src/gpt/liturgy/roman-rite/postconciliar/roman-missal-third-edition-en-us-2011/propers/general-calendar/shared/formularies/pc-r07-all-saints/propers/verified.md sha256=3025d42355a3162576f155c12ed1d0eadca32ec592546ef8ca6bf7c7066e5630
% appointment-dependency: src/gpt/liturgy/roman-rite/postconciliar/roman-missal-third-edition-en-us-2011/propers/registry/README.md sha256=4e2d92326dece103e1475e70ffb5e447815f9b12d697926e48a54850a1641208
% appointment-dependency: src/gpt/liturgy/roman-rite/postconciliar/roman-missal-third-edition-en-us-2011/propers/registry/formula-dispositions.md sha256=de58754f5a8ab61dffa3f4a6a479c843fe3763b6a5b37099282960de71d93989
% appointment-note: entrance-antiphon: Composed Missal text without an appointed biblical citation; theological echoes and compiler chant directions do not create appointments.
% appointment-note: collect: Composed Missal text without an appointed biblical citation; theological echoes and compiler chant directions do not create appointments.
% appointment-note: first-reading: Revelation 7:2-4,9-14; appointed; vulgate numbering; canonical verse queries, not appointed wording.
% appointment-note: responsorial-psalm: Psalm 24:1bc-2,3-4ab,5-6; appointed; hebrew numbering; whole-verse query envelope, not subverse chronology or appointed wording.
% appointment-note: responsorial-psalm: Cf. Psalm 24:6; adaptation; hebrew numbering; canonical verse queries, not appointed wording.
% appointment-note: second-reading: 1 John 3:1-3; appointed; vulgate numbering; canonical verse queries, not appointed wording.
% appointment-note: gospel-acclamation: Matthew 11:28; appointed; vulgate numbering; canonical verse queries, not appointed wording.
% appointment-note: gospel: Matthew 5:1-12a; appointed; vulgate numbering; whole-verse query envelope, not subverse chronology or appointed wording.
% appointment-note: prayer-over-offerings: Composed Missal text without an appointed biblical citation; theological echoes and compiler chant directions do not create appointments.
% appointment-note: proper-preface: Composed Missal text without an appointed biblical citation; theological echoes and compiler chant directions do not create appointments.
% appointment-note: communion-antiphon: Matthew 5:8-10; appointed; vulgate numbering; canonical verse queries, not appointed wording.
% appointment-note: prayer-after-communion: Composed Missal text without an appointed biblical citation; theological echoes and compiler chant directions do not create appointments.
% appointment-note: solemn-blessing: Composed Missal text without an appointed biblical citation; theological echoes and compiler chant directions do not create appointments.
% first-reading: First Reading
\expandafter\def\csname triptychchronologyannotation@first-reading\endcsname{%
  \chronologyannotationgroup{composition}{disputed}{\textbf{Composition} -- disputed: \chronologyannotationreach{Apoc.7.10}{Apoc}{true}\chronologyannotationreach{Apoc.7.11}{Apoc}{true}\chronologyannotationreach{Apoc.7.12}{Apoc}{true}\chronologyannotationreach{Apoc.7.13}{Apoc}{true}\chronologyannotationreach{Apoc.7.14}{Apoc}{true}\chronologyannotationreach{Apoc.7.2}{Apoc}{true}\chronologyannotationreach{Apoc.7.3}{Apoc}{true}\chronologyannotationreach{Apoc.7.4}{Apoc}{true}\chronologyannotationreach{Apoc.7.9}{Apoc}{true}\chronologyannotationclaim{composition.apocalypse}{composition}{catholic-traditional-v1}{disputed}{the reign of Domitian ... A.D. 81-96}{A.D. 81--96}; \chronologyannotationreach{Apoc.7.10}{Apoc}{true}\chronologyannotationreach{Apoc.7.11}{Apoc}{true}\chronologyannotationreach{Apoc.7.12}{Apoc}{true}\chronologyannotationreach{Apoc.7.13}{Apoc}{true}\chronologyannotationreach{Apoc.7.14}{Apoc}{true}\chronologyannotationreach{Apoc.7.2}{Apoc}{true}\chronologyannotationreach{Apoc.7.3}{Apoc}{true}\chronologyannotationreach{Apoc.7.4}{Apoc}{true}\chronologyannotationreach{Apoc.7.9}{Apoc}{true}\chronologyannotationclaim{composition.apocalypse}{composition}{catholic-traditional-v1}{disputed}{from the year 90 to 100 (approximately)}{c. A.D. 90--100}.}%
}
% responsorial-psalm: Responsorial Psalm
\expandafter\def\csname triptychchronologyannotation@responsorial-psalm\endcsname{%
  \chronologyannotationgroup{historical-setting}{research-pending}{\textbf{Historical setting} -- Occasion date unresolved.}%
  \space \chronologyannotationgroup{composition}{preferred}{\textbf{Composition}: \chronologyannotationreach{Ps.23.1}{Ps}{true}\chronologyannotationreach{Ps.23.2}{Ps}{true}\chronologyannotationreach{Ps.23.3}{Ps}{true}\chronologyannotationreach{Ps.23.4}{Ps}{true}\chronologyannotationreach{Ps.23.5}{Ps}{true}\chronologyannotationreach{Ps.23.6}{Ps}{true}\chronologyannotationclaim{critical.psalms.latest-composition-boundary}{composition}{catholic-critical-v1}{preferred}{before the Maccabean period, around 165 B.C.; no individual psalm can be dated securely}{Before c. 165 B.C.}}%
}
% second-reading: Second Reading
\expandafter\def\csname triptychchronologyannotation@second-reading\endcsname{%
  \chronologyannotationgroup{composition}{preferred}{\textbf{Composition}: \chronologyannotationreach{1John.3.1}{1John}{true}\chronologyannotationreach{1John.3.2}{1John}{true}\chronologyannotationreach{1John.3.3}{1John}{true}\chronologyannotationclaim{composition.first-epistle-of-st-john}{composition}{catholic-traditional-v1}{preferred}{from the year 90 to 100 (approximately)}{c. A.D. 90--100}.}%
}
% gospel-acclamation: Gospel acclamation
\expandafter\def\csname triptychchronologyannotation@gospel-acclamation\endcsname{%
  \chronologyannotationgroup{composition}{disputed}{\textbf{Composition} -- disputed: \chronologyannotationreach{Matt.11.28}{Matt}{true}\chronologyannotationclaim{composition.gospel-of-matthew}{composition}{catholic-traditional-v1}{disputed}{about A.D. 38-45}{c. A.D. 38--45}; \chronologyannotationreach{Matt.11.28}{Matt}{true}\chronologyannotationclaim{composition.gospel-of-matthew}{composition}{catholic-traditional-v1}{disputed}{about the year 40-42}{c. A.D. 40--42}; \chronologyannotationreach{Matt.11.28}{Matt}{true}\chronologyannotationclaim{composition.gospel-of-matthew}{composition}{catholic-traditional-v1}{disputed}{the years 40-45}{A.D. 40--45}; \chronologyannotationreach{Matt.11.28}{Matt}{true}\chronologyannotationclaim{composition.gospel-of-matthew}{composition}{catholic-traditional-v1}{disputed}{about the year 60-68}{c. A.D. 60--68}; \chronologyannotationreach{Matt.11.28}{Matt}{true}\chronologyannotationclaim{composition.gospel-of-matthew}{composition}{catholic-traditional-v1}{disputed}{about the years 64-67}{c. A.D. 64--67}; \chronologyannotationreach{Matt.11.28}{Matt}{true}\chronologyannotationclaim{composition.gospel-of-matthew}{composition}{catholic-traditional-v1}{disputed}{about the year 50}{c. A.D. 50}.}%
}
% gospel: Gospel
\expandafter\def\csname triptychchronologyannotation@gospel\endcsname{%
  \chronologyannotationgroup{narrated-event}{research-pending}{\textbf{Event} -- No narrated-event date in the chronology corpus.}%
  \space \chronologyannotationgroup{composition}{disputed}{\textbf{Composition} -- disputed: \chronologyannotationreach{Matt.5.1}{Matt}{true}\chronologyannotationreach{Matt.5.10}{Matt}{true}\chronologyannotationreach{Matt.5.11}{Matt}{true}\chronologyannotationreach{Matt.5.12}{Matt}{true}\chronologyannotationreach{Matt.5.2}{Matt}{true}\chronologyannotationreach{Matt.5.3}{Matt}{true}\chronologyannotationreach{Matt.5.4}{Matt}{true}\chronologyannotationreach{Matt.5.5}{Matt}{true}\chronologyannotationreach{Matt.5.6}{Matt}{true}\chronologyannotationreach{Matt.5.7}{Matt}{true}\chronologyannotationreach{Matt.5.8}{Matt}{true}\chronologyannotationreach{Matt.5.9}{Matt}{true}\chronologyannotationclaim{composition.gospel-of-matthew}{composition}{catholic-traditional-v1}{disputed}{about A.D. 38-45}{c. A.D. 38--45}; \chronologyannotationreach{Matt.5.1}{Matt}{true}\chronologyannotationreach{Matt.5.10}{Matt}{true}\chronologyannotationreach{Matt.5.11}{Matt}{true}\chronologyannotationreach{Matt.5.12}{Matt}{true}\chronologyannotationreach{Matt.5.2}{Matt}{true}\chronologyannotationreach{Matt.5.3}{Matt}{true}\chronologyannotationreach{Matt.5.4}{Matt}{true}\chronologyannotationreach{Matt.5.5}{Matt}{true}\chronologyannotationreach{Matt.5.6}{Matt}{true}\chronologyannotationreach{Matt.5.7}{Matt}{true}\chronologyannotationreach{Matt.5.8}{Matt}{true}\chronologyannotationreach{Matt.5.9}{Matt}{true}\chronologyannotationclaim{composition.gospel-of-matthew}{composition}{catholic-traditional-v1}{disputed}{about the year 40-42}{c. A.D. 40--42}; \chronologyannotationreach{Matt.5.1}{Matt}{true}\chronologyannotationreach{Matt.5.10}{Matt}{true}\chronologyannotationreach{Matt.5.11}{Matt}{true}\chronologyannotationreach{Matt.5.12}{Matt}{true}\chronologyannotationreach{Matt.5.2}{Matt}{true}\chronologyannotationreach{Matt.5.3}{Matt}{true}\chronologyannotationreach{Matt.5.4}{Matt}{true}\chronologyannotationreach{Matt.5.5}{Matt}{true}\chronologyannotationreach{Matt.5.6}{Matt}{true}\chronologyannotationreach{Matt.5.7}{Matt}{true}\chronologyannotationreach{Matt.5.8}{Matt}{true}\chronologyannotationreach{Matt.5.9}{Matt}{true}\chronologyannotationclaim{composition.gospel-of-matthew}{composition}{catholic-traditional-v1}{disputed}{the years 40-45}{A.D. 40--45}; \chronologyannotationreach{Matt.5.1}{Matt}{true}\chronologyannotationreach{Matt.5.10}{Matt}{true}\chronologyannotationreach{Matt.5.11}{Matt}{true}\chronologyannotationreach{Matt.5.12}{Matt}{true}\chronologyannotationreach{Matt.5.2}{Matt}{true}\chronologyannotationreach{Matt.5.3}{Matt}{true}\chronologyannotationreach{Matt.5.4}{Matt}{true}\chronologyannotationreach{Matt.5.5}{Matt}{true}\chronologyannotationreach{Matt.5.6}{Matt}{true}\chronologyannotationreach{Matt.5.7}{Matt}{true}\chronologyannotationreach{Matt.5.8}{Matt}{true}\chronologyannotationreach{Matt.5.9}{Matt}{true}\chronologyannotationclaim{composition.gospel-of-matthew}{composition}{catholic-traditional-v1}{disputed}{about the year 60-68}{c. A.D. 60--68}; \chronologyannotationreach{Matt.5.1}{Matt}{true}\chronologyannotationreach{Matt.5.10}{Matt}{true}\chronologyannotationreach{Matt.5.11}{Matt}{true}\chronologyannotationreach{Matt.5.12}{Matt}{true}\chronologyannotationreach{Matt.5.2}{Matt}{true}\chronologyannotationreach{Matt.5.3}{Matt}{true}\chronologyannotationreach{Matt.5.4}{Matt}{true}\chronologyannotationreach{Matt.5.5}{Matt}{true}\chronologyannotationreach{Matt.5.6}{Matt}{true}\chronologyannotationreach{Matt.5.7}{Matt}{true}\chronologyannotationreach{Matt.5.8}{Matt}{true}\chronologyannotationreach{Matt.5.9}{Matt}{true}\chronologyannotationclaim{composition.gospel-of-matthew}{composition}{catholic-traditional-v1}{disputed}{about the years 64-67}{c. A.D. 64--67}; \chronologyannotationreach{Matt.5.1}{Matt}{true}\chronologyannotationreach{Matt.5.10}{Matt}{true}\chronologyannotationreach{Matt.5.11}{Matt}{true}\chronologyannotationreach{Matt.5.12}{Matt}{true}\chronologyannotationreach{Matt.5.2}{Matt}{true}\chronologyannotationreach{Matt.5.3}{Matt}{true}\chronologyannotationreach{Matt.5.4}{Matt}{true}\chronologyannotationreach{Matt.5.5}{Matt}{true}\chronologyannotationreach{Matt.5.6}{Matt}{true}\chronologyannotationreach{Matt.5.7}{Matt}{true}\chronologyannotationreach{Matt.5.8}{Matt}{true}\chronologyannotationreach{Matt.5.9}{Matt}{true}\chronologyannotationclaim{composition.gospel-of-matthew}{composition}{catholic-traditional-v1}{disputed}{about the year 50}{c. A.D. 50}.}%
}
% communion-antiphon: Communion antiphon
\expandafter\def\csname triptychchronologyannotation@communion-antiphon\endcsname{%
  \chronologyannotationgroup{composition}{disputed}{\textbf{Composition} -- disputed: \chronologyannotationreach{Matt.5.10}{Matt}{true}\chronologyannotationreach{Matt.5.8}{Matt}{true}\chronologyannotationreach{Matt.5.9}{Matt}{true}\chronologyannotationclaim{composition.gospel-of-matthew}{composition}{catholic-traditional-v1}{disputed}{about A.D. 38-45}{c. A.D. 38--45}; \chronologyannotationreach{Matt.5.10}{Matt}{true}\chronologyannotationreach{Matt.5.8}{Matt}{true}\chronologyannotationreach{Matt.5.9}{Matt}{true}\chronologyannotationclaim{composition.gospel-of-matthew}{composition}{catholic-traditional-v1}{disputed}{about the year 40-42}{c. A.D. 40--42}; \chronologyannotationreach{Matt.5.10}{Matt}{true}\chronologyannotationreach{Matt.5.8}{Matt}{true}\chronologyannotationreach{Matt.5.9}{Matt}{true}\chronologyannotationclaim{composition.gospel-of-matthew}{composition}{catholic-traditional-v1}{disputed}{the years 40-45}{A.D. 40--45}; \chronologyannotationreach{Matt.5.10}{Matt}{true}\chronologyannotationreach{Matt.5.8}{Matt}{true}\chronologyannotationreach{Matt.5.9}{Matt}{true}\chronologyannotationclaim{composition.gospel-of-matthew}{composition}{catholic-traditional-v1}{disputed}{about the year 60-68}{c. A.D. 60--68}; \chronologyannotationreach{Matt.5.10}{Matt}{true}\chronologyannotationreach{Matt.5.8}{Matt}{true}\chronologyannotationreach{Matt.5.9}{Matt}{true}\chronologyannotationclaim{composition.gospel-of-matthew}{composition}{catholic-traditional-v1}{disputed}{about the years 64-67}{c. A.D. 64--67}; \chronologyannotationreach{Matt.5.10}{Matt}{true}\chronologyannotationreach{Matt.5.8}{Matt}{true}\chronologyannotationreach{Matt.5.9}{Matt}{true}\chronologyannotationclaim{composition.gospel-of-matthew}{composition}{catholic-traditional-v1}{disputed}{about the year 50}{c. A.D. 50}.}%
}
% presentation: orientation-v2
\newcommand{\chronologyorientationrow}[2]{%
  \ifcsname triptychchronologyorientation@#1\endcsname
    \csname triptychchronologyorientation@#1\endcsname{#2}%
  \else
    \PackageError{triptych}{No chronology orientation row for #1}{}%
  \fi}
\expandafter\def\csname triptychchronologyorientation@row-f5cbb3bd9936e1a8\endcsname#1{%
Responsorial Psalm: Psalm 24:1bc-2,3--4ab,5--6, Cf. Psalm 24:6 (Hebrew numbering; Vulgate 23) & Historical setting. #1 & Occasion date unresolved\\[.25em]%
}
\expandafter\def\csname triptychchronologyorientation@row-e3f82d7d7c816aeb\endcsname#1{%
Responsorial Psalm: Psalm 24:1bc-2,3--4ab,5--6, Cf. Psalm 24:6 (Hebrew numbering; Vulgate 23) & Composition: Psalms. #1 & \chronologyannotationreach{Ps.23.1}{Ps}{true}\chronologyannotationreach{Ps.23.2}{Ps}{true}\chronologyannotationreach{Ps.23.3}{Ps}{true}\chronologyannotationreach{Ps.23.4}{Ps}{true}\chronologyannotationreach{Ps.23.5}{Ps}{true}\chronologyannotationreach{Ps.23.6}{Ps}{true}\chronologyannotationclaim{critical.psalms.latest-composition-boundary}{composition}{catholic-critical-v1}{preferred}{before the Maccabean period, around 165 B.C.; no individual psalm can be dated securely}{Before c. 165 B.C.}\\[.25em]%
}
\expandafter\def\csname triptychchronologyorientation@row-5a9c87b5fd78fefb\endcsname#1{%
Gospel: Matthew 5:1--12a & Event. #1 & Date unresolved\\[.25em]%
}
\expandafter\def\csname triptychchronologyorientation@row-6661a9700ce1b710\endcsname#1{%
Gospel acclamation: Matthew 11:28; Gospel: Matthew 5:1--12a; Communion antiphon: Matthew 5:8--10 & Composition: The Gospel of St Matthew. #1 & Disputed: \chronologyannotationreach{Matt.11.28}{Matt}{true}\chronologyannotationreach{Matt.5.1}{Matt}{true}\chronologyannotationreach{Matt.5.10}{Matt}{true}\chronologyannotationreach{Matt.5.11}{Matt}{true}\chronologyannotationreach{Matt.5.12}{Matt}{true}\chronologyannotationreach{Matt.5.2}{Matt}{true}\chronologyannotationreach{Matt.5.3}{Matt}{true}\chronologyannotationreach{Matt.5.4}{Matt}{true}\chronologyannotationreach{Matt.5.5}{Matt}{true}\chronologyannotationreach{Matt.5.6}{Matt}{true}\chronologyannotationreach{Matt.5.7}{Matt}{true}\chronologyannotationreach{Matt.5.8}{Matt}{true}\chronologyannotationreach{Matt.5.9}{Matt}{true}\chronologyannotationclaim{composition.gospel-of-matthew}{composition}{catholic-traditional-v1}{disputed}{about A.D. 38-45}{c. A.D. 38--45}; \chronologyannotationreach{Matt.11.28}{Matt}{true}\chronologyannotationreach{Matt.5.1}{Matt}{true}\chronologyannotationreach{Matt.5.10}{Matt}{true}\chronologyannotationreach{Matt.5.11}{Matt}{true}\chronologyannotationreach{Matt.5.12}{Matt}{true}\chronologyannotationreach{Matt.5.2}{Matt}{true}\chronologyannotationreach{Matt.5.3}{Matt}{true}\chronologyannotationreach{Matt.5.4}{Matt}{true}\chronologyannotationreach{Matt.5.5}{Matt}{true}\chronologyannotationreach{Matt.5.6}{Matt}{true}\chronologyannotationreach{Matt.5.7}{Matt}{true}\chronologyannotationreach{Matt.5.8}{Matt}{true}\chronologyannotationreach{Matt.5.9}{Matt}{true}\chronologyannotationclaim{composition.gospel-of-matthew}{composition}{catholic-traditional-v1}{disputed}{about the year 40-42}{c. A.D. 40--42}; \chronologyannotationreach{Matt.11.28}{Matt}{true}\chronologyannotationreach{Matt.5.1}{Matt}{true}\chronologyannotationreach{Matt.5.10}{Matt}{true}\chronologyannotationreach{Matt.5.11}{Matt}{true}\chronologyannotationreach{Matt.5.12}{Matt}{true}\chronologyannotationreach{Matt.5.2}{Matt}{true}\chronologyannotationreach{Matt.5.3}{Matt}{true}\chronologyannotationreach{Matt.5.4}{Matt}{true}\chronologyannotationreach{Matt.5.5}{Matt}{true}\chronologyannotationreach{Matt.5.6}{Matt}{true}\chronologyannotationreach{Matt.5.7}{Matt}{true}\chronologyannotationreach{Matt.5.8}{Matt}{true}\chronologyannotationreach{Matt.5.9}{Matt}{true}\chronologyannotationclaim{composition.gospel-of-matthew}{composition}{catholic-traditional-v1}{disputed}{the years 40-45}{A.D. 40--45}; \chronologyannotationreach{Matt.11.28}{Matt}{true}\chronologyannotationreach{Matt.5.1}{Matt}{true}\chronologyannotationreach{Matt.5.10}{Matt}{true}\chronologyannotationreach{Matt.5.11}{Matt}{true}\chronologyannotationreach{Matt.5.12}{Matt}{true}\chronologyannotationreach{Matt.5.2}{Matt}{true}\chronologyannotationreach{Matt.5.3}{Matt}{true}\chronologyannotationreach{Matt.5.4}{Matt}{true}\chronologyannotationreach{Matt.5.5}{Matt}{true}\chronologyannotationreach{Matt.5.6}{Matt}{true}\chronologyannotationreach{Matt.5.7}{Matt}{true}\chronologyannotationreach{Matt.5.8}{Matt}{true}\chronologyannotationreach{Matt.5.9}{Matt}{true}\chronologyannotationclaim{composition.gospel-of-matthew}{composition}{catholic-traditional-v1}{disputed}{about the year 60-68}{c. A.D. 60--68}; \chronologyannotationreach{Matt.11.28}{Matt}{true}\chronologyannotationreach{Matt.5.1}{Matt}{true}\chronologyannotationreach{Matt.5.10}{Matt}{true}\chronologyannotationreach{Matt.5.11}{Matt}{true}\chronologyannotationreach{Matt.5.12}{Matt}{true}\chronologyannotationreach{Matt.5.2}{Matt}{true}\chronologyannotationreach{Matt.5.3}{Matt}{true}\chronologyannotationreach{Matt.5.4}{Matt}{true}\chronologyannotationreach{Matt.5.5}{Matt}{true}\chronologyannotationreach{Matt.5.6}{Matt}{true}\chronologyannotationreach{Matt.5.7}{Matt}{true}\chronologyannotationreach{Matt.5.8}{Matt}{true}\chronologyannotationreach{Matt.5.9}{Matt}{true}\chronologyannotationclaim{composition.gospel-of-matthew}{composition}{catholic-traditional-v1}{disputed}{about the years 64-67}{c. A.D. 64--67}; \chronologyannotationreach{Matt.11.28}{Matt}{true}\chronologyannotationreach{Matt.5.1}{Matt}{true}\chronologyannotationreach{Matt.5.10}{Matt}{true}\chronologyannotationreach{Matt.5.11}{Matt}{true}\chronologyannotationreach{Matt.5.12}{Matt}{true}\chronologyannotationreach{Matt.5.2}{Matt}{true}\chronologyannotationreach{Matt.5.3}{Matt}{true}\chronologyannotationreach{Matt.5.4}{Matt}{true}\chronologyannotationreach{Matt.5.5}{Matt}{true}\chronologyannotationreach{Matt.5.6}{Matt}{true}\chronologyannotationreach{Matt.5.7}{Matt}{true}\chronologyannotationreach{Matt.5.8}{Matt}{true}\chronologyannotationreach{Matt.5.9}{Matt}{true}\chronologyannotationclaim{composition.gospel-of-matthew}{composition}{catholic-traditional-v1}{disputed}{about the year 50}{c. A.D. 50}\\[.25em]%
}
\expandafter\def\csname triptychchronologyorientation@row-4f310c075dd990dc\endcsname#1{%
Second Reading: 1 John 3:1--3 & Composition: The First Epistle of St John. #1 & \chronologyannotationreach{1John.3.1}{1John}{true}\chronologyannotationreach{1John.3.2}{1John}{true}\chronologyannotationreach{1John.3.3}{1John}{true}\chronologyannotationclaim{composition.first-epistle-of-st-john}{composition}{catholic-traditional-v1}{preferred}{from the year 90 to 100 (approximately)}{c. A.D. 90--100}\\[.25em]%
}
\expandafter\def\csname triptychchronologyorientation@row-02f156d681e4660a\endcsname#1{%
First Reading: Revelation 7:2--4,9--14 & Composition: The Apocalypse of St John. #1 & Disputed: \chronologyannotationreach{Apoc.7.10}{Apoc}{true}\chronologyannotationreach{Apoc.7.11}{Apoc}{true}\chronologyannotationreach{Apoc.7.12}{Apoc}{true}\chronologyannotationreach{Apoc.7.13}{Apoc}{true}\chronologyannotationreach{Apoc.7.14}{Apoc}{true}\chronologyannotationreach{Apoc.7.2}{Apoc}{true}\chronologyannotationreach{Apoc.7.3}{Apoc}{true}\chronologyannotationreach{Apoc.7.4}{Apoc}{true}\chronologyannotationreach{Apoc.7.9}{Apoc}{true}\chronologyannotationclaim{composition.apocalypse}{composition}{catholic-traditional-v1}{disputed}{the reign of Domitian ... A.D. 81-96}{A.D. 81--96}; \chronologyannotationreach{Apoc.7.10}{Apoc}{true}\chronologyannotationreach{Apoc.7.11}{Apoc}{true}\chronologyannotationreach{Apoc.7.12}{Apoc}{true}\chronologyannotationreach{Apoc.7.13}{Apoc}{true}\chronologyannotationreach{Apoc.7.14}{Apoc}{true}\chronologyannotationreach{Apoc.7.2}{Apoc}{true}\chronologyannotationreach{Apoc.7.3}{Apoc}{true}\chronologyannotationreach{Apoc.7.4}{Apoc}{true}\chronologyannotationreach{Apoc.7.9}{Apoc}{true}\chronologyannotationclaim{composition.apocalypse}{composition}{catholic-traditional-v1}{disputed}{from the year 90 to 100 (approximately)}{c. A.D. 90--100}\\[.25em]%
}

Composition, narrated event and remembered setting answer different historical questions. The table locates each biblical appointment once, consolidating overlapping uses in Matthew. The following notes explain the source judgments and alternatives behind the displayed dates.
\begin{chronologyorientationtable}
\chronologyorientationrow{row-f5cbb3bd9936e1a8}{Seeking admission to the Lord's mountain and sanctuary; original occasion unknown. See Psalm note below.}
\chronologyorientationrow{row-e3f82d7d7c816aeb}{Individual writing place unknown; the date is a bound for the collection. See Psalm note below.}
\chronologyorientationrow{row-5a9c87b5fd78fefb}{Jesus teaches disciples on a mountain during his Galilean ministry; exact hill unknown. See Matthew note below.}
\chronologyorientationrow{row-6661a9700ce1b710}{Tradition associates an earlier Hebrew form with Palestine; the Greek text's relation is disputed. See Matthew note below.}
\chronologyorientationrow{row-4f310c075dd990dc}{Probably associated with Ephesus; the immediate recipients' town is uncertain. See First John note below.}
\chronologyorientationrow{row-02f156d681e4660a}{The seer names Patmos; the book addresses seven churches in Asia. See Revelation note below.}
\end{chronologyorientationtable}

\subsection{Psalm 24: a sanctuary setting and a collection bound}
\label{history-psalm}
The Douay--Rheims superscription names David and the first day of the week. Augustine reads that first day through the Lord's resurrection and glorification. That is his Christian exposition of the superscription, and does not identify a dated occasion of the original psalm. The ascent, holy place and seekers describe worship at God's sanctuary. An exact occasion and individual writing place remain unknown.\footnote{Psalm 23, superscription and verses 1--6; Augustine, \work{Expositions}, Psalm 23.1--6.}

The collection boundary above derives from the official Psalms introduction's account of the collection before the Maccabean period and its warning that individual psalms cannot be dated securely. It is a latest bound, not the year this psalm was composed. The traditional author label and that collection boundary therefore describe different facts. No reign-era date has been substituted for a missing occasion.\footnote{USCCB, NABRE, ``The Book of Psalms,'' introduction, discussion of dating and collection; exact public delivery inspected for this research on 26 October 2026.}

\subsection{Matthew: the mountain and the writing of the Gospel}
\label{history-matthew}
The event row concerns Jesus' teaching on a mountain. Matthew 4:23--5:2 places it in the Galilean ministry and names the disciples who approach him, with the crowds in context. It names no particular hill. The event's exact year is unresolved in the chronology evidence. The scene is a narrated occurrence; it is separate from the question when the evangelist wrote the Gospel.

The composition alternatives preserve the differing ancient and modern reckonings reported by Jacquier. His discussion includes the tradition of a Palestinian Hebrew original and the disputed relation of the Greek Gospel to it. The earlier and later estimates cannot all be treated as agreed dates of one demonstrably identical textual event. No particular one is preferred here. The table retains their disagreement rather than silently assigning a proposed Hebrew date to the extant Greek work.\footnote{E. Jacquier, ``Gospel of St. Matthew,'' \work{Catholic Encyclopedia}, X (1911), discussions of original language, date and place of composition.}

The acclamation from chapter 11 and Communion's repetition of three blessings share the Gospel's composition question. They create no independent composition occasions. Chrysostom's exposition of the acclamation reads it with the Son's revelation of the Father and the following invitation to learn his humility. That literary setting helps explain the verse, while adding no exact date to the event table.

\subsection{First John: probable Ephesian setting}
\label{history-john}
The date displayed derives from Durand's approximate span for the Johannine writings collectively. It is not a separately pinpointed occasion for this epistle. Drum considers an Ephesian association and the end of the first century probable, rather than certain. The immediate recipients' town is not established. First John 3:1--3 is instruction about present sonship and future manifestation; it narrates no separately datable earthly event.\footnote{A. Durand, ``The New Testament,'' \work{Catholic Encyclopedia}, XIV (1912), ``Origin,'' on the Johannine writings; Walter Drum, ``Epistles of Saint John,'' VIII (1910), First Epistle, ``Time and Place'' and ``Destination and Purpose.''}

\subsection{Revelation: Patmos and the churches in Asia}
\label{history-apocalypse}
Revelation 1:9 identifies Patmos in the seer's account; the opening addresses seven churches in Asia. The throne in chapter 7 belongs to the vision. It is neither a recovered writing room nor a date for the saints' arrival in heaven. The heavenly scene and the circumstances in which the book was written remain distinct.

Van den Biesen reports the traditional Domitianic window while discussing the difficulty raised by a Vespasian interpretation. Durand gives the approximate Johannine span displayed as another composition alternative. Their overlap is compatible with a late first-century setting, but the two sources have different scopes and do not provide one agreed exact year. The traditional attribution to John and historical questions about composition are preserved as the cited sources state them.\footnote{F. Van den Biesen, ``Apocalypse,'' \work{Catholic Encyclopedia}, I (1907), date and authorship discussion; Durand, ``The New Testament,'' ``Origin.''}
